\documentclass[11pt]{article}
\usepackage[margin=1in]{geometry}
\usepackage{amsmath,amssymb,amsthm,mathtools}
\usepackage{booktabs}
\usepackage{microtype}
\usepackage[hidelinks]{hyperref}
\newtheorem{lemma}{Lemma}
\newtheorem{proposition}[lemma]{Proposition}
\newcommand{\eps}{\varepsilon}
\newcommand{\lmax}{\lambda_{\max}}
\newcommand{\Ran}{\operatorname{Ran}}

\title{Does a block-spin map preserve the dynamics?\\[2pt]
\large A two-spacing test with SU(3) glueball operators}
\author{[Authors]\\ \small AlphaOmega Institute}
\date{Draft, \today}

\begin{document}
\maketitle

\begin{abstract}
\noindent
Blocking fixes which coarse-lattice operator a fine-lattice operator represents. If the map preserves
the dynamics, evolving a coarse state and then translating it, or translating it first and evolving the
fine state over the same physical time, should agree. We estimate the largest disagreement over a
window of operators from correlators of two independently simulated lattices, through an identity that
is exact for vacuum transfer moments, with the operator correspondence fixed by blocking and aligned by
its polar decomposition. For SU(3) Yang--Mills theory with the Wilson action at $\beta=5.7$ and $6.054$,
in the primary window of two smeared 0$^{++}$ operators, both eigenvalues of the blocked fine lattice
lie below their coarse partners, so a comparison of spectra shows only faster fine decay. Yet the
blocking-fixed pairing rotates the leading direction by $24.8(1.0)^\circ$, and in one paired direction the
normalized one-step autocorrelation is $0.146(5)$ on the fine lattice against $0.101(6)$ on the coarse
one, a one-sided discrepancy of $0.045(6)$. The fitted scalar masses, which show a known finite-spacing
shift, reproduce the size of the leading spectral separation within uncertainty. Different blocking
recipes rank in opposite order by the faithfulness of the initial pairing and by the tracking of the
evolution, and again by the symmetric and the one-sided discrepancy. We state a sufficient one-sided
criterion for a gap, show what correlators determine and what they cannot (an upper bound on slow
states the chosen operators do not see), and record the design and predictions for a second pair,
$\beta=6.054$ against $6.546$.
\end{abstract}

\section{Introduction}

Lattice computations of glueball masses compare numbers between lattice spacings and extrapolate. The
maps between the lattices, which renormalization-group arguments presuppose, are rarely tested as maps.
A spatial blocking step is such a map: it fixes, before any data are seen, which fine-lattice operator
represents a given coarse one. It can then be asked whether the map preserves what happens to the
states it pairs. Evolve a coarse state over one coarse time step and translate it to the fine lattice;
or translate it first and evolve the fine state over the same physical interval. The distance between
the two results, maximized over the states of a window of operators, measures how far the map fails to
preserve the dynamics. A second maximization, over all coarse states with the map extended by zero
outside the window, also counts states that start outside the coarse window and evolve into it; it is
the defect $\|S_fW-WS_c\|$ of the window map. We report both: on the primary recipe of the first pair
they are $0.180(11)$ and $0.203(13)$.

The defect is not determined by the two spectra. Our clearest example is the first lattice pair studied
here. In the window spanned by two smeared 0$^{++}$ operators, both eigenvalues of the blocked fine
lattice's one-step transfer compression lie below their coarse partners ($0.2360$ and $0.0459$ against
$0.3392$ and $0.0636$): paired in spectral order, the fine lattice decays faster in every direction. Paired
as blocking prescribes, the leading direction is rotated by $24.8(1.0)^\circ$, and in one direction the fine
evolution decays more slowly than the coarse one: the blocking-paired preparation that maximizes the
signed discrepancy has normalized one-step autocorrelation $0.101(6)$ on the coarse lattice and $0.146(5)$
on the fine one. This coexists with a heavier fitted fine glueball, because the paired preparation is not
the ground state of either lattice. A comparison of masses, or of spectra, cannot see it.

The question has been seen before. Monte Carlo renormalization-group matching between two lattice
spacings was developed in the 1980s; glueball correlations were considered as matching observables but
set aside because their long-distance signal is statistically hard, and blocked loop expectation values
were matched instead \cite{Mackenzie1984,Swendsen1981}. The aim was to tune the coarse action, and no
blocking kernel was expected to yield exactly the Wilson action at finite spacing, since that would make
the Wilson action free of finite-spacing errors \cite{Mackenzie1984}. The two-lattice method has been
revisited for SU(3) gauge theories \cite{Hasenfratz2009}. Effective transfer matrices on block variables
have been estimated by Monte Carlo on one lattice \cite{HPR1994}. With symmetric normalization, our
one-lattice leakage matrix is the Gram matrix of the first block residual in the block Lanczos method of
Hackett and Wagman \cite{HackettWagman2025}. What we add is the combination: a correspondence fixed by
blocking, aligned by its polar part, and compared through correlators measured independently at two
spacings, giving the defect of that map, including evolution out of either window, and the part of the
disagreement that is visible inside the paired windows.

The defect is also the hypothesis of renormalization-group arguments for a mass gap, in which a gap is
carried along a chain of maps between ever finer lattices. Only one direction of disagreement matters
for a lower bound on a gap, and we give the corresponding one-sided criterion, its operator form for a
chain, and the quantity that simulation cannot supply. Nothing here bears directly on the continuum
theory; the contribution is a measurement and the precise statement of what it tests.

Section~\ref{sec:framework} sets out the framework, Section~\ref{sec:setup} the lattice setup,
Section~\ref{sec:pair1} the results for $\beta=5.7$ against $6.054$, and Section~\ref{sec:pair2} the
design and predictions, written before its fine data were examined, for $\beta=6.054$ against $6.546$.

\section{Framework}\label{sec:framework}

\subsection{Operator windows}

Let $O_1,\dots,O_n$ be zero-momentum operators on time slices, centered by subtracting their vacuum
expectation values, with correlator matrix $C_{kl}(t)=\langle O_k(t)O_l(0)\rangle-\langle O_k\rangle\langle
O_l\rangle$. With the transfer matrix $T$ normalized to unit vacuum eigenvalue and vacuum dominance
assumed, $C(t)=\Phi^*T^t\Phi$, where the columns of $\Phi$ are the states the operators create; they are
orthogonal to the vacuum. With $G=C(0)$ and the orthonormal frame $U=\Phi G^{-1/2}$, the window
compression and leakage at lag $\ell$ are
\begin{align}
 K(\ell)&=G^{-1/2}C(\ell)G^{-1/2}=U^*T^\ell U,\\
 L(\ell)&=G^{-1/2}C(2\ell)G^{-1/2}-K(\ell)^2=U^*T^\ell(I-UU^*)T^\ell U\ \ge 0 .
\end{align}
$L$ measures how far one step of evolution carries the window out of itself. These identities hold for
vacuum transfer moments; the measured $K$ and $L$ are finite-sample estimates from periodic correlators,
and negative eigenvalues of an estimated $L$ (noise) are set to zero before square roots are taken. At
finite periodic time extent the ensemble correlators are thermal traces, and the vacuum reading is an
approximation; for a single level the backward term alone gives $\sqrt{L}\approx10^{-3}$ at the lattices
used here.

\subsection{The correspondence fixed by blocking}

A blocking map $B$ takes a fine configuration to a coarse one, so a coarse operator $O$ defines the fine
operator $O\circ B$. Measured on the fine lattice, these operators give a correlator matrix $C_f$ and Gram
matrix $G_f$; the coarse lattice gives $C_c$ and $G_c$. In orthonormal frames the correspondence of
operator states is $A=G_f^{1/2}G_c^{-1/2}$, with polar decomposition $A=VM$, $V$ unitary and $M\ge0$. The
window map is $W=U_fVU_c^*$. With spacings $a_c=2a_f$, coarse lag $\ell$ and fine lag $2\ell$ cover the same
physical interval; write $S_c=T_c^{\ell}$ and $S_f=T_f^{2\ell}$, and $K_c,L_c,K_f,L_f$ for the corresponding
compressions and leakages.

\subsection{The defect and the spectral floor}

\begin{proposition}\label{prop:defect}
With $D=K_fV-VK_c$,
\begin{equation}
 \|S_fW-WS_c\|=\left\|\begin{pmatrix} D & -V\sqrt{L_c}\\ \sqrt{L_f}\,V & 0\end{pmatrix}\right\|
 \le \|D\|+\max\bigl(\|L_f\|,\|L_c\|\bigr)^{1/2}.
\end{equation}
\end{proposition}

Here $W$ is extended by zero outside the coarse window, so the norm runs over all coarse inputs, and a
state outside the window that evolves into it contributes through the block $-V\sqrt{L_c}$. Restricted
to inputs in the coarse window the disagreement is
\begin{equation}
 \|(S_fW-WS_c)U_c\|=\lmax\bigl(D^*D+V^*L_fV\bigr)^{1/2},
\end{equation}
which is at most the defect. The proof (Appendix~\ref{app:proofs}) uses the polar decompositions of the
leakage maps $(I-U_iU_i^*)S_iU_i$. Three quantities have distinct meanings, though their norms do not add. The minimum
of $\|V'^*K_fV'-K_c\|$ over all unitary $V'$ is the largest difference between matched sorted eigenvalues
of $K_f$ and $K_c$ \cite{Bhatia}, a floor set by the spectra. The alignment $V$ selected by blocking gives
the in-window mismatch $\|K_fV-VK_c\|$. Evolution out of the windows enters the defect through $L_c,L_f$.
The angle between the leading eigenvectors of $K_c$ and $V^*K_fV$ measures how far the blocking-fixed
pairing departs from a pairing in spectral order.

In the doubled space $\mathcal H_c\oplus\mathcal H_f$ with $S=S_c\oplus S_f$, the polar reflection
$R=\begin{psmallmatrix}0&W^*\\W&0\end{psmallmatrix}$ has the aligned and anti-aligned graphs
$\{(x,\pm Wx)\}$ as eigenspaces, and $\|[S,R]\|=\|S_fW-WS_c\|$: the defect measures how far evolution fails
to commute with the operation that exchanges each matched coarse description with its fine partner. The
normalization $\|M-I\|$ measures how far the raw correspondence is from isometric, which is the
near-neutrality of its graph for the form $\operatorname{diag}(I,-I)$; the defect measures the
near-invariance of the polar graph (Appendix~\ref{app:doubled}).

\subsection{The one-sided criterion}

For a lower bound on a gap only one direction matters: whether the fine evolution could decay more
slowly than the coarse one. Let $0\le T\le I$ be a transfer operator over the matched interval with
$T\Omega=\Omega$, let $P$ be a window orthogonal to $\Omega$ and $Q$ the rest of $\Omega^\perp$.

\begin{lemma}\label{lem:onesided}
With $k=\lmax(PTP)$, $\ell^2=\|QTP\|^2$ and $q=\|QTQ\|$, the top of $T$ on $\Omega^\perp$ obeys
$\lambda_2(T)\le\mu=\tfrac{k+q}{2}+\sqrt{(\tfrac{k-q}{2})^2+\ell^2}$, and $\mu\le s$ exactly when $k\le s$,
$q\le s$ and $\ell^2\le(s-k)(s-q)$. Positivity of $T$ implies $L\le qK$ in window coordinates, so
$q\ge\lmax(K^{-1/2}LK^{-1/2})$ when $K>0$ (with the inverse taken on the support of $K$ otherwise).
\end{lemma}

The data bound $q$ from below; a gap certificate needs it bounded from above. When the fine window is
inferred from the coarse one, only the positive part of $D'=V^*K_fV-K_c$ enters: $V^*K_fV=K_c+D'\le
K_c+D'_+$, so $k\le\lmax(K_c+D'_+)\le\lmax(K_c)+\eps_+$ with $\eps_+=\max(0,\lmax(D'))$. The matrix form
keeps the direction of the error.

\begin{proposition}\label{prop:chain}
On the vacuum's complement, with $a_c=2a_f$, suppose $T_c\le B_c$, $V^*U_f^*T_f^2U_fV-U_c^*T_cU_c\le E$,
$L=U_f^*T_f^2QT_f^2U_f$ and $QT_f^2Q\le q$. Then for every $\gamma>0$
\begin{equation}
 T_f^2\le\bigl(C+L/\gamma\bigr)\oplus(q+\gamma)I_Q,\qquad C=V(U_c^*B_cU_c+E)V^*,
\end{equation}
so that $T_f\le B_f=(C+L/\gamma)^{1/2}\oplus\sqrt{q+\gamma}\,I_Q$. If $s=e^{-a_cm_*}>q$ and $C+L/(s-q)\le s$,
then $T_f\le e^{-a_fm_*}$.
\end{proposition}

This is a sufficient criterion. $B_f$ is a full operator bound, so it can be carried to the next rung of a
chain. A uniform gap needs one $m_*$ through every refinement, volume and sector claimed, and an
independent upper bound on $q$ at every rung. Carrying the directional information across a shared
lattice requires the same operator family on both sides of it: where one pair ends on blocked operators
and the next starts on unblocked ones, as at $\beta=6.054$ below, the cross-Gram matrix between the two
families on the same configurations is needed. In a chain the inequality must run from coarse to fine.
In recent work organized around block-spin maps and summable defects \cite{FaizalShabir2026}, the
one-step estimate (their Section~5, Eqs.~(5.35)--(5.36)) lower-bounds the coarser gap using the finer
gap; by itself it cannot propagate a known coarse lower bound to finer lattices. Identifications between different Hilbert
spaces that control spectra are the subject of quasi-unitary equivalence \cite{Post2012,PostSimmer};
$\|M-I\|$ and the defect play the roles of near-unitarity and approximate intertwining there.

\section{Lattice setup}\label{sec:setup}

We simulate SU(3) Yang--Mills theory with the Wilson plaquette action \cite{Wilson1974}, each sweep one
heatbath pass \cite{Creutz1980,CabibboMarinari1982,KennedyPendleton1985} and three overrelaxation passes
\cite{BrownWoch1987}. The 0$^{++}$ operators are slice sums of the real traces of spatial plaquettes after
$n$ steps of spatial stout smearing \cite{MorningstarPeardon2004} with $\rho=0.1$, at $n=4,8,16,24$.

Blocking acts in the spatial directions only. A coarse spatial link is the SU(3) projection of the
straight two-link path plus $\alpha$ times the four bent paths of length four between the same sites,
$\alpha=0,0.3,0.6$ \cite{Teper1987}. The projection is the polar part with its determinant phase removed,
which is covariant at both ends of a path; when the path sum is numerically singular, or when the
determinant of the polar part lies within $10^{-8}$ of $-1$ in phase (where the cube root jumps by a
center element), the straight path is used instead. No such fallback occurred in any run. Blocking with
each of the eight origins of the coarse sublattice and averaging the operator amplitudes gives the
zero-momentum projection; the single origin gives the fixed-origin variant, which folds in momentum
$\pi/a_f$.

\begin{table}[t]
\centering
\begin{tabular}{@{}lllll@{}}
\toprule
$\beta$ & lattice & configurations & sweeps apart & role \\
\midrule
5.7 & $8^3\times16$ & 3000 & 4 & first pair, coarse \\
6.054 & $16^3\times32$, blocked to $8^3$ & 3000 & 4 & first pair, fine \\
6.054 & $16^3\times32$ & 3000 & 4 & second pair, coarse \\
6.546 & $32^3\times64$, blocked to $16^3$ & 1000 & 8 & second pair, fine (in progress) \\
\bottomrule
\end{tabular}
\caption{Ensembles. Spacing ratios from the Necco--Sommer interpolation \cite{NeccoSommer2002}:
$a(5.7)/a(6.054)=2.0015$, $a(6.054)/a(6.546)=1.9997$; all boxes have $L\sqrt\sigma\approx3.2$.}
\label{tab:ensembles}
\end{table}

Table~\ref{tab:ensembles} lists the ensembles. Errors come from a delete-one jackknife over bins of ten
configurations, applied to both ensembles of a pair with the variances added; all blocking recipes are
computed from the same fine bins and are resampled together, which gives paired errors for their
differences. Masses are fitted to $A\cosh(m(t-T/2))$ with a covariance from the jackknife, the
variational vector fixed on the full ensemble \cite{Michael1985,LuscherWolff1990,Blossier2009}. The data,
code and analysis outputs are public (Appendix~\ref{app:data}); rerun from the published data on a second
machine, the analyses reproduce the outputs of the first byte for byte.

\section{Results: \texorpdfstring{$\beta=5.7$}{beta = 5.7} against \texorpdfstring{$\beta=6.054$}{beta = 6.054}}\label{sec:pair1}

\subsection{Masses}

The coarse lattice gives $am_0=0.973(35)$ at $t=1$ and a correlated fit over $t\in[1,4]$ of $0.969$ in
coarse units, consistent with $0.987(9)$ measured by Athenodorou and Teper at $\beta=5.6924$
\cite{AthenodorouTeper2020}. The blocked fine lattice has a plateau at $a_fm_0\approx0.66$
($0.674(13)$, $0.658(24)$, $0.663(49)$ at $t=1,2,3$ for $\alpha=0.6$ with zero momentum), consistent with
about $0.645$, an interpolation to $\beta=6.054$ of published values ($0.6365(43)$ at $\beta=6.0625$
\cite{AthenodorouTeper2020}). Over matched physical intervals the masses
differ (Table~\ref{tab:masses}): a fit over fine $t\in[2,8]$ against coarse $t\in[1,4]$ gives a
difference of $0.171(29)$ in fine units, $m_f/m_c-1=0.354(68)$ with 90\% interval $[0.242,0.466]$; with the
first interval removed (fine $[4,8]$ against coarse $[2,4]$) the central value persists at lower
significance, $0.150(79)$. In units of the string tension the coarse scalar is near $m_0/\sqrt\sigma\approx2.4$
and the fine one near $3.3$: the scalar dip of the Wilson action near $\beta\approx5.7$, where the scalar
softens selectively \cite{BhanotCreutz1981,Heller1995,AthenodorouTeper2020}. Paired within each ensemble,
the leading window eigenvalues differ by $0.1032(65)$ and the single-state decay factors of the fitted
masses over the matched interval by $0.1100(181)$; observed minus predicted is $-0.0068(163)$. The
individual window eigenvalues lie below their single-state predictions by $0.0404(120)$ (coarse) and
$0.0336(110)$ (fine), offsets that largely cancel in the difference. The fitted scalar masses reproduce the
size of the leading separation within uncertainty: quantitative support for the scalar-dip reading,
which leaves its causal share open. We attribute nothing beyond the leading separation to it.

\begin{table}[t]
\centering
\begin{tabular}{@{}ccccc@{}}
\toprule
interval $j$ & fine, $\alpha=0.6$ zero & coarse & difference & difference, $\alpha=0.3$ zero \\
\midrule
0 & 0.722(9) & 0.540(7) & 0.182(12) & 0.198(12) \\
1 & 0.660(32) & 0.486(17) & 0.174(37) & 0.178(37) \\
2 & 0.730(134) & 0.457(44) & 0.27(14) & 0.24(14) \\
\bottomrule
\end{tabular}
\caption{Effective masses per fine step over matched intervals (fine $t=2j\to2j+2$, coarse
$t=j\to j+1$), in fine units.}
\label{tab:masses}
\end{table}

\subsection{Window quantities}

\begin{table}[t]
\centering
\small\setlength{\tabcolsep}{4.5pt}
\begin{tabular}{@{}lccccccc@{}}
\toprule
$\alpha$, blocking & $\|M-I\|$ & $\|K_fV-VK_c\|$ & floor & window inputs & defect & departure & $\eps_+$ \\
\midrule
0, fixed & 1.725(11) & 0.301(6) & 0.294 & 0.314(7) & 0.333(9) & 0.183(7) & 0 \\
0, zero & 0.288(6) & 0.252(6) & 0.196 & 0.268(8) & 0.286(10) & 0.165(8) & 0.010(6) \\
0.3, fixed & 0.273(5) & 0.208(6) & 0.172 & 0.239(8) & 0.256(10) & 0.156(9) & 0 \\
0.3, zero & 0.180(4) & 0.173(6) & 0.110 & 0.187(10) & 0.209(12) & 0.129(11) & 0.040(6) \\
0.6, fixed & 0.141(4) & 0.193(6) & 0.150 & 0.225(9) & 0.243(10) & 0.150(9) & 0.008(6) \\
0.6, zero & 0.246(4) & 0.166(6) & 0.103 & 0.180(11) & 0.203(13) & 0.126(11) & 0.045(6) \\
\bottomrule
\end{tabular}
\caption{Window $\{4,8\}$ at one coarse step. ``Floor'': the minimum mismatch over alignments;
``window inputs'': the disagreement maximized over coarse-window states; ``defect'': maximized over all
coarse states; ``departure'': $\|(I-P_+)SP_+\|$ for the aligned graph; $\eps_+$: the positive part of
$\lmax(V^*K_fV-K_c)$. The coarse window compared with itself gives a departure of $0.122$.}
\label{tab:window}
\end{table}

Table~\ref{tab:window} collects the window quantities. The in-window mismatch, $0.17$--$0.30$, is far
above the noise reference: pairs of independent resamples of 3000 configurations of the coarse ensemble,
whose true mismatch is zero, give an apparent mismatch with 95th percentile $0.017$. The window
$\{4,8,16,24\}$ gives the same picture (mismatch $0.166(6)$, disagreement over window inputs $0.177(11)$,
defect $0.197(13)$). The two maximizations differ by about $0.02$ for every recipe: states outside the
coarse window that evolve into it add little to what the window states themselves show.

\subsection{What the blocking map shows beyond the spectra}

For $\alpha=0.6$ with zero momentum, both fine eigenvalues lie below the coarse ones, and pairing the
eigenvectors in spectral order gives $\eps_+=0$ exactly. The blocking-fixed pairing gives
$\eps_+=0.045(6)$, because the leading direction of $V^*K_fV$ is rotated by $24.8(1.0)^\circ$ from that of
$K_c$ ($23.0(1.1)^\circ$ for the fixed origin, $25.2(1.0)^\circ$ for $\alpha=0.3$). Concretely, in the
direction that maximizes the signed discrepancy the blocking-paired preparation has normalized one-step
autocorrelation $0.101(6)$ on the coarse lattice and $0.146(5)$ on the fine one. That preparation mixes
the leading and subleading window states differently on the two lattices, so its fine autocorrelation
can exceed the coarse one although both spectra point the other way. At twice the matched
interval the angle stays at $23.0(3.0)^\circ$ while $\eps_+$ falls to $0.0092(62)$. Evolving twice inside
each measured window, with $K_f^2$ and $K_c^2$ and the same alignment, predicts $0.0127(13)$; the measured
value, which also allows evolution to leave the windows and return, differs by $-0.0035(60)$. The window
modes, evolving at their own rates, account for the decline within errors.

\subsection{Recipes rank in opposite order}

Averaging over origins lowers the mismatch, the defect and the departure at every $\alpha$; at
$\alpha=0.6$ the paired differences are $0.0273(25)$, $0.0396(53)$ and $0.0243(45)$. The fixed origin
nevertheless gives the more nearly isometric raw correspondence ($\|M-I\|=0.141$ against $0.246$): judged
by how faithfully it pairs the states at the initial instant, the map would be chosen wrongly for
tracking them one step later. The one-sided error ranks the recipes the other way again:
$\eps_+(\text{fixed})-\eps_+(\text{zero})=-0.0369(27)$ at $\alpha=0.6$ and $-0.0401(57)$ at $\alpha=0.3$.
A gap certificate depends on $\eps_+$ and the leakage together, and the fixed origin has twice the fine
leakage, so which recipe serves a certificate better is not settled by these errors.

\subsection{The one-sided certificate}

For the reference $s=e^{-0.973}$, the coarse mass over the matched interval, Lemma~\ref{lem:onesided} with
the fine window measured directly needs $q\le0.30$--$0.35$ depending on the recipe ($0.326(27)$ for
$\alpha=0.6$ with zero momentum), above the positivity floors ($0.05$--$0.24$, noisy). Inferred from the
coarse lattice through the scalar bound $\lmax(K_c)+\eps_+=0.384>s$, no $q$ suffices; the matrix form
gives $\lmax(K_c+D'_+)=0.3464(45)$ and, with the window test of Proposition~\ref{prop:chain}, the
conditional ceiling $q\le0.167(117)$. Part of the apparent loss of information in passing from the coarse
lattice came from replacing a matrix by its norm.

\subsection{Checks}

Bins of 20 and 50 change no central value; the headline errors grow by at most about 8\%, and the
largest growth in the full output is 23\%. The four-operator window reproduces the reversal of $\eps_+$
($0.056(6)$ with zero momentum against $0.014(6)$ and $0.007(5)$ with the fixed origin); its estimated
leakage matrices have negative eigenvalues within errors, which the estimator sets to zero.

\section{Second pair: \texorpdfstring{$\beta=6.054$}{beta = 6.054} against \texorpdfstring{$\beta=6.546$}{beta = 6.546}}\label{sec:pair2}

The first pair's coarse lattice lies in the scalar dip. The second pair places both lattices in the
scaling region: a new $\beta=6.054$ ensemble on $16^3\times32$, measured unblocked, against $\beta=6.546$
on $32^3\times64$ blocked to $16^3\times64$ with $\alpha=0.3,0.6$. The following were recorded before the
fine data were examined.

\emph{Prediction.} Published masses give $m_0/\sqrt\sigma\approx3.27$ at $\beta=6.0625$ and, extrapolated,
$\approx3.36$ at $6.546$ \cite{AthenodorouTeper2020}, so single-state decay factors over one coarse step
differ by about $0.003$ (first pair $0.110$). If the scalar dip drives the first pair's leading separation,
the leading window eigenvalues should agree within about $0.03$ (first pair $0.103$), and the scalar masses
within 5\%. No prediction is made for the angle, $\eps_+$ or the leakage.

\emph{Endpoints.} Primary comparison: $\alpha=0.6$ with zero momentum, window $\{4,8\}$, one coarse step.
Mass endpoint $r=m_f/m_c-1$ from fits over fine $[4,16]$ against coarse $[2,8]$, with delayed-start, even-time
and short-range checks; a two-sided 90\% interval wholly inside $[-0.05,0.05]$ supports agreement, wholly
outside it a discrepancy, partial overlap is unresolved. The new pair is also compared at the first pair's
physical interval (coarse lag 2, fine lag 4), with the reference mass $0.655$ per $\beta=6.054$ step.
Caveats: published scalar masses shift by about 4\% between $L\sqrt\sigma\approx2.7$ and $3.9$
\cite{AthenodorouTeper2020}, so a 5\% tolerance is not calibrated at our volume; and the spacing ratio
is the value of an interpolation whose uncertainty we do not propagate \cite{NeccoSommer2002}, so the
predicted $0.003$ is illustrative and the mass verdict is conditional on the ratio.

\emph{Results.} [To be completed after the recorded analysis.]

\section{Discussion}

The measurement answers a question that masses do not: whether a fixed map between two lattice
descriptions preserves the evolution of the states it pairs. On the first pair the answer is no, for two
reasons that the method separates. The leading spectral separation is a known finite-spacing effect in
the scalar channel. The rotation of the leading direction under the blocking-fixed pairing, and the
one-sided discrepancy it produces, are properties of the map. The second pair tests whether the first
shrinks when both lattices scale, and whether the second shrinks with it.

For renormalization-group arguments the results sharpen three points. The inequality that carries a gap
must run from coarse to fine. Errors should be carried as operator inequalities, not norms, since a norm
charges a harmless direction as if it were the worst one. And the hypothesis that no correlator can
supply is an upper bound on how slowly states outside the chosen windows can evolve. A slow eigenstate
orthogonal to every window state is invisible to every window correlator: appending a decoupled
eigenstate with transfer eigenvalue arbitrarily close to one leaves every measured correlator unchanged
while the full gap shrinks. (An arbitrary orthogonal vector, by contrast, can evolve back into the
window.) What an independent argument would have to establish is a quantitative lower bound on how much
of every low-energy state, including linear combinations, the blocked operators capture.

Limitations: one channel, finite and modest volumes, the Wilson action, finite-time windows of two or
four operators, and statistical errors; the scale ratios come from an interpolation.

\section{Conclusion}

A blocking-fixed correspondence between two lattices, aligned by its polar decomposition, can be tested
exactly through correlators. On $\beta=5.7$ against $6.054$ it reveals a rotation of the leading window
direction and a one-sided discrepancy that the spectra hide, beside a spectral separation explained by the
scalar dip. The one-sided criterion states what such measurements can establish for a gap, and what they
cannot.

\appendix

\section{Proofs}\label{app:proofs}

\emph{Proposition~\ref{prop:defect}.} Write $Y_i\sqrt{L_i}$ for the polar decomposition of
$(I-U_iU_i^*)S_iU_i$, with $Y_i$ a partial isometry into the complement of the window that acts
isometrically on the support of $L_i$. Then $S_fW-WS_c=U_f(K_fV-VK_c)U_c^*+Y_f\sqrt{L_f}VU_c^*-
U_fV\sqrt{L_c}Y_c^*$. In the systems $(U_f,Y_f)$ and $(U_c,Y_c)$ this is the block matrix shown; it
annihilates the part of the input space outside the support of $\sqrt{L_c}$ and has range in the
corresponding output support, so the norms are equal also when the leakage matrices are singular. The
bound follows from the triangle inequality, and the window-input formula from restricting to inputs
$U_ca$.

\emph{Spectral floor.} For Hermitian $A,B$ and unitary $V'$, $\|V'^*AV'-B\|\ge\max_i|\lambda_i^\downarrow(A)-
\lambda_i^\downarrow(B)|$, with equality when $V'$ maps eigenvectors in order \cite{Bhatia}.

\emph{Lemma~\ref{lem:onesided}.} For a unit $\psi=u+v\perp\Omega$ with $u=P\psi$, $v=Q\psi$,
$\langle\psi,T\psi\rangle\le k\|u\|^2+2\ell\|u\|\|v\|+q\|v\|^2$, at most the top eigenvalue of
$\begin{psmallmatrix}k&\ell\\ \ell&q\end{psmallmatrix}$. Positivity of $T$ makes its block form positive, so
$QTP=(QTQ)^{1/2}X(PTP)^{1/2}$ with $\|X\|\le1$, whence $(QTP)^*(QTP)\le q\,PTP$.

\emph{Proposition~\ref{prop:chain}.} Compressing $T_c\le B_c$ and adding $E$ bounds the window block of
$T_f^2$ by $C$. For $\gamma>0$, $\begin{psmallmatrix}X&Y^*\\Y&Z\end{psmallmatrix}\le
\begin{psmallmatrix}X+Y^*Y/\gamma&0\\0&Z+\gamma\end{psmallmatrix}$, the difference being
$\begin{psmallmatrix}\gamma^{-1/2}Y^*\\-\gamma^{1/2}\end{psmallmatrix}
\begin{psmallmatrix}\gamma^{-1/2}Y&-\gamma^{1/2}\end{psmallmatrix}\ge0$; the operator square root is
monotone, and $\gamma=s-q$ gives the window test.

\section{The doubled space}\label{app:doubled}

[To be written from the review guide: the polar reflection and its graphs, the identity
$\|[S,R]\|=\|S_fW-WS_c\|$, the Krein square of the raw graph, and the transfer of the aligned graph into
the anti-aligned graph ($\tfrac12\|K_c-V^*K_fV\|$) and out of the windows
($\lmax(\tfrac12(L_c+V^*L_fV))^{1/2}$).]

\section{Data, code and reproduction}\label{app:data}

Code, data and analysis outputs are in the repository\\
\url{https://github.com/AlphaOmega-Institute/AlphaOmega-LQCD}\\
(\texttt{lattice/}, \texttt{lattice/results/}). [Version and archive identifiers to be added.]

\emph{Use of AI assistance.} [Statement to be completed according to the venue's policy.]

\begin{thebibliography}{99}
\bibitem{Wilson1974} K.~G. Wilson, Phys.\ Rev.\ D \textbf{10} (1974) 2445.
\bibitem{Creutz1980} M.~Creutz, Phys.\ Rev.\ D \textbf{21} (1980) 2308.
\bibitem{CabibboMarinari1982} N.~Cabibbo and E.~Marinari, Phys.\ Lett.\ B \textbf{119} (1982) 387.
\bibitem{KennedyPendleton1985} A.~D. Kennedy and B.~J. Pendleton, Phys.\ Lett.\ B \textbf{156} (1985) 393.
\bibitem{BrownWoch1987} F.~R. Brown and T.~J. Woch, Phys.\ Rev.\ Lett.\ \textbf{58} (1987) 2394.
\bibitem{MorningstarPeardon2004} C.~Morningstar and M.~Peardon, Phys.\ Rev.\ D \textbf{69} (2004) 054501.
\bibitem{Teper1987} M.~Teper, Phys.\ Lett.\ B \textbf{183} (1987) 345.
\bibitem{Michael1985} C.~Michael, Nucl.\ Phys.\ B \textbf{259} (1985) 58.
\bibitem{LuscherWolff1990} M.~L\"uscher and U.~Wolff, Nucl.\ Phys.\ B \textbf{339} (1990) 222.
\bibitem{Blossier2009} B.~Blossier, M.~Della Morte, G.~von Hippel, T.~Mendes and R.~Sommer, JHEP \textbf{04} (2009) 094.
\bibitem{NeccoSommer2002} S.~Necco and R.~Sommer, Nucl.\ Phys.\ B \textbf{622} (2002) 328.
\bibitem{AthenodorouTeper2020} A.~Athenodorou and M.~Teper, JHEP \textbf{11} (2020) 172, arXiv:2007.06422.
\bibitem{BhanotCreutz1981} G.~Bhanot and M.~Creutz, Phys.\ Rev.\ D \textbf{24} (1981) 3212.
\bibitem{Heller1995} U.~M. Heller, arXiv:hep-lat/9508009.
\bibitem{Mackenzie1984} P.~B. Mackenzie, Monte Carlo renormalization group calculations for SU(2) lattice gauge theory, FERMILAB-Conf-84/48-T (1984).
\bibitem{Swendsen1981} R.~H. Swendsen, Phys.\ Rev.\ Lett.\ \textbf{47} (1981) 1775.
\bibitem{Hasenfratz2009} A.~Hasenfratz, arXiv:0907.0919.
\bibitem{HPR1994} M.~Hasenbusch, K.~Pinn and K.~Rummukainen, Phys.\ Rev.\ D \textbf{49} (1994) 495.
\bibitem{HackettWagman2025} D.~C. Hackett and M.~L. Wagman, Phys.\ Rev.\ D \textbf{112} (2025) 014514, arXiv:2412.04444.
\bibitem{FaizalShabir2026} M.~Faizal and A.~Shabir, Fortschr.\ Phys.\ \textbf{74} (2026) e70097, arXiv:2606.19362.
\bibitem{Post2012} O.~Post, \emph{Spectral Analysis on Graph-like Spaces}, Lecture Notes in Math.\ 2039, Springer (2012).
\bibitem{PostSimmer} O.~Post and J.~Simmer, arXiv:2503.22611.
\bibitem{Bhatia} R.~Bhatia, \emph{Matrix Analysis}, Springer (1997).
\end{thebibliography}

\end{document}
